% !TeX root = ../教材题库_全部条目_转写工作稿.tex
% Source: XJ-B1-C04-RV-P151-Y1430
% 第4章 / 章末复习 复习题四 / 教材 p.151、152 / PDF 第158、159页
% OCR未校对：公式与文字均可能有误；未完成数学TeX转写；禁止直接发给学生
% 原题及相邻原解截图（核对用）：../题目原图/XJ-B1-C04-RV-P151-Y1430-01.webp, ../题目原图/XJ-B1-C04-RV-P151-Y1430-02.webp
% 本文件只在转写完成后作为题干使用；OCR 条目尚不是可靠题干。
\begin{bankitem}{复习题四 第9题}{第4章 / 章末复习 复习题四 / 教材 p.151、152 / PDF 第158、159页}{XJ-B1-C04-RV-P151-Y1430}{OCR未校对：公式与文字均可能有误；未完成数学TeX转写；禁止直接发给学生}
\begin{ocrdraft}
[教材p.151 / PDF第158页]
\par
9．如图，有一条曲线是函数3，一a'（0<
\par
0
\par
a<D的图象，其他三条曲线是从这条曲线
\par
出发经轴反射得到的．试写出这些曲线对应
\par
的函数表达式．
\par
（第9题）
\par
[教材p.152 / PDF第159页]
\end{ocrdraft}
\end{bankitem}
